% Options for packages loaded elsewhere
\PassOptionsToPackage{unicode}{hyperref}
\PassOptionsToPackage{hyphens}{url}
%
\documentclass[
]{article}
\usepackage{amsmath,amssymb}
\usepackage{iftex}
\ifPDFTeX
  \usepackage[T1]{fontenc}
  \usepackage[utf8]{inputenc}
  \usepackage{textcomp} % provide euro and other symbols
\else % if luatex or xetex
  \usepackage{unicode-math} % this also loads fontspec
  \defaultfontfeatures{Scale=MatchLowercase}
  \defaultfontfeatures[\rmfamily]{Ligatures=TeX,Scale=1}
\fi
\usepackage{lmodern}
\ifPDFTeX\else
  % xetex/luatex font selection
\fi
% Use upquote if available, for straight quotes in verbatim environments
\IfFileExists{upquote.sty}{\usepackage{upquote}}{}
\IfFileExists{microtype.sty}{% use microtype if available
  \usepackage[]{microtype}
  \UseMicrotypeSet[protrusion]{basicmath} % disable protrusion for tt fonts
}{}
\makeatletter
\@ifundefined{KOMAClassName}{% if non-KOMA class
  \IfFileExists{parskip.sty}{%
    \usepackage{parskip}
  }{% else
    \setlength{\parindent}{0pt}
    \setlength{\parskip}{6pt plus 2pt minus 1pt}}
}{% if KOMA class
  \KOMAoptions{parskip=half}}
\makeatother
\usepackage{xcolor}
\setlength{\emergencystretch}{3em} % prevent overfull lines
\providecommand{\tightlist}{%
  \setlength{\itemsep}{0pt}\setlength{\parskip}{0pt}}
\setcounter{secnumdepth}{-\maxdimen} % remove section numbering
\ifLuaTeX
  \usepackage{selnolig}  % disable illegal ligatures
\fi
\IfFileExists{bookmark.sty}{\usepackage{bookmark}}{\usepackage{hyperref}}
\IfFileExists{xurl.sty}{\usepackage{xurl}}{} % add URL line breaks if available
\urlstyle{same}
\hypersetup{
  hidelinks,
  pdfcreator={LaTeX via pandoc}}

\author{}
\date{}

\begin{document}

---\ntitle: ``Linear Algebra for CV, DL, and AI''\nauthor:
Santhosh\ninstitute: LA1000\n---\n\n\# Week 3: Matrix Transformations
(Slides)

\hypertarget{slide-1-introduction-to-matrix-transformations}{%
\subsection{Slide 1: Introduction to Matrix
Transformations}\label{slide-1-introduction-to-matrix-transformations}}

\begin{itemize}
\tightlist
\item
  Matrices as transformations in space:

  \begin{itemize}
  \tightlist
  \item
    Reflections
  \item
    Rotations
  \item
    Scaling
  \item
    Shearing
  \end{itemize}
\item
  Reference Video:
  \href{https://www.youtube.com/watch?v=xeyFqlHw8C0}{Gilbert Strang -
  Matrix Multiplication}
\end{itemize}

\begin{center}\rule{0.5\linewidth}{0.5pt}\end{center}

\hypertarget{slide-2-reflection}{%
\subsection{Slide 2: Reflection}\label{slide-2-reflection}}

\begin{itemize}
\tightlist
\item
  Definition: Matrix that flips vectors across a line/plane.
\item
  Example:

  \begin{itemize}
  \tightlist
  \item
    Across x-axis: (

    \begin{bmatrix} 1 & 0 \\ 0 & -1 \end{bmatrix}

    )
  \item
    Visualization Tool:
    \href{../labs/supplementary_notebooks/week3_matrix_transformations.ipynb}{Week
    3 Notebook}.
  \end{itemize}
\end{itemize}

\begin{center}\rule{0.5\linewidth}{0.5pt}\end{center}

\hypertarget{slide-3-rotation}{%
\subsection{Slide 3: Rotation}\label{slide-3-rotation}}

\begin{itemize}
\tightlist
\item
  Rotational Matrices rotate vectors counterclockwise by theta: (

  \begin{bmatrix} \cos{\theta} & -\sin{\theta} \\ \sin{\theta} & \cos{\theta} \end{bmatrix}

  )
\item
  Applications:

  \begin{itemize}
  \tightlist
  \item
    Image processing rotations.
  \item
    Computer Graphics.
  \end{itemize}
\end{itemize}

\begin{center}\rule{0.5\linewidth}{0.5pt}\end{center}

\hypertarget{slide-4-scaling}{%
\subsection{Slide 4: Scaling}\label{slide-4-scaling}}

\begin{itemize}
\tightlist
\item
  Definition: Stretch or compress vectors.
\item
  Applications:

  \begin{itemize}
  \tightlist
  \item
    Adjusting shapes in graphical designs.
  \item
    Uniform vs Non-uniform scaling.
  \end{itemize}
\end{itemize}

\begin{center}\rule{0.5\linewidth}{0.5pt}\end{center}

\hypertarget{slide-5-summary-exercises}{%
\subsection{Slide 5: Summary \&
Exercises}\label{slide-5-summary-exercises}}

\begin{enumerate}
\def\labelenumi{\arabic{enumi}.}
\tightlist
\item
  Practice transformations using provided code cells.
\item
  Solve eigenvalue problems to solidify understanding.
\end{enumerate}

\end{document}
